\section{Extended-Precision Floating-Point Types for Perturbation}
\label{sec:extended-precision-types}

In perturbation-based rendering (\cref{sec:perturbation-concept}), each pixel
computes a low-precision \emph{delta} from a high-precision reference orbit.
The precision of this delta arithmetic (the \code{SubType} template parameter)
determines both the accuracy of the rendered image and the GPU throughput.
\FractalShark{} provides a hierarchy of multi-component floating-point types
that trade additional ALU operations for higher effective precision without
relying on native 64-bit floating-point hardware.

\subsection{Motivation: the perturbation precision gap}
\label{subsec:ep-motivation}

Single-precision IEEE~754 (\code{float}, 24-bit mantissa) is the fastest
arithmetic available on consumer GPUs, but it does not always provide enough
precision for accurate perturbation rendering.  In some cases,
rendering artifacts appear because the per-pixel deltas lose significant bits
during the perturbation recurrence.  Native \code{double} (53-bit mantissa)
resolves these artifacts but incurs a substantial throughput penalty on consumer
GPUs, where fp64 performance is significantly slower than fp32, with throughput ratios on the order of 32:1 on some architectures.

The \textbf{2$\times$32 double-float} type bridges this gap: it provides
${\sim}48$~bits of effective mantissa using only fp32 arithmetic, delivering
sufficient precision for accurate perturbation while remaining on the fast
fp32 ALUs.  This format is the extended-precision type with real practical value in
\FractalShark{}'s rendering pipeline.

\FractalShark{} also includes higher-precision types (2$\times$64,
4$\times$32, 4$\times$64) as reference implementations based on existing work,
primarily for experimentation and completeness.  Some base (non-perturbation)
kernels also use these types for direct Mandelbrot iteration (e.g.\
\code{mandel\_2x\_float}, \code{mandel\_4x\_float}), but those kernels are not
particularly useful in practice---perturbation with linear approximation is the
primary rendering path at any serious zoom depth.

\subsection{The 2\texorpdfstring{$\times$}{x}32 double-float type (\code{dblflt})}
\label{subsec:ep-dblflt}

The core type is \code{MattDblflt} (aliased as \code{dblflt}), a struct
containing two \code{float} values:
\begin{description}
  \item[\code{head}] The most significant bits of the represented value.
  \item[\code{tail}] A correction term satisfying
    $|\mathtt{tail}| \le \tfrac{1}{2}\,\mathrm{ulp}(\mathtt{head})$,
    so that the represented value is exactly $\mathtt{head} + \mathtt{tail}$.
\end{description}
Construction from a \code{double} uses Knuth's error-free sum
algorithm \cite{Knuth1997TAOCP2} to
split the value into a head--tail pair.  The total storage is 8~bytes---the
same footprint as a single \code{double}.

\paragraph{Arithmetic operations.}
The CUDA device functions in \code{dblflt.cuh} implement the standard
operations:

\begin{center}
\begin{tabular}{lll}
\hline
\textbf{Function} & \textbf{Operation} & \textbf{Max relative error} \\
\hline
\code{add\_dblflt(a, b)} & $a + b$  & $O(u^2) \approx 2^{-48}$ \\
\code{sub\_dblflt(a, b)} & $a - b$  & $O(u^2) \approx 2^{-48}$ \\
\code{mul\_dblflt(a, b)} & $a \times b$  & $O(u^2) \approx 2^{-47}$ \\
\code{sqr\_dblflt(a)}    & $a^2$         & optimized via $2 \cdot \mathtt{head} \cdot \mathtt{tail}$ \\
\code{mul\_dblflt2x(a,b)}& $2ab$         & multiply then shift \\
\hline
\end{tabular}
\end{center}

All operations are built on fused multiply-add (\code{\_\_fmaf\_rn},
\code{\_\_fmaf\_rd}) intrinsics, which compute $a \times b + c$ in a single
rounding step.  This property is essential for maintaining the head--tail invariant:
intermediate products must be exact or correctly rounded before the correction
term is computed.  The error bounds listed above are the theoretical
$O(u^2)$ bounds for compensated float arithmetic, where $u = 2^{-24}$ is
the float unit roundoff.  Note that the source code comments in
\code{dblflt.cuh} cite empirical error values of ${\sim}2^{-104.7}$;
however, these measurements appear to have been inherited from the
double-double (\code{dbldbl}) implementation (where they are consistent
with $u = 2^{-53}$) and were not re-measured for the float-float variant.

\subsection{\code{CudaDblflt} wrapper and \code{HDRFloat} integration}
\label{subsec:ep-cudadblflt}

The \code{CudaDblflt<T>} class template wraps a \code{dblflt} value with
C++ operator overloads (\code{+}, \code{-}, \code{*}, \code{/}, comparisons)
suitable for use in GPU device code.  This wrapper allows perturbation kernels to be
written generically over the \code{SubType} parameter without
\code{dblflt}-specific function calls.

For perturbation rendering with linear approximation, the delta values need
both extended precision \emph{and} wide dynamic range.  This combination is achieved by
combining \code{CudaDblflt} with the HDR wrapper
(\cref{sec:hdr-float}):
\begingroup\footnotesize
\[
  \code{HDRFloat<CudaDblflt<dblflt>>} \;=\;
    \underbrace{\text{2$\times$32-bit mantissa}}_{\sim\text{48~bits}}
    \;+\;
    \underbrace{\text{int32 exponent}}_{\text{wide dynamic range}}
    \;=\; \text{12~bytes total.}
\]
\endgroup
The \code{Reduce()} method on \code{HDRFloat} normalizes the two float
components relative to the shared exponent by extracting and aligning their
IEEE~754 exponent fields.  A \code{DoubleTo2x32Converter} utility handles
host-side fallback paths where CUDA device code is unavailable.

In the render algorithm enumeration, kernel variants using this type are named
with the \code{2x32} or \code{HDRx2x32} suffix (e.g.\
\code{Gpu2x32PerturbedLAv2}, \code{GpuHDRx2x32PerturbedLAv2}).

\subsection{Higher-precision types (academic interest)}
\label{subsec:ep-higher-prec}

\FractalShark{} includes additional multi-component types based on existing
implementations.  These are primarily of academic interest and are not
essential for typical rendering.

\paragraph{2\texorpdfstring{$\times$}{x}64 double-double (\code{dbldbl}).}
\code{MattDbldbl} pairs two \code{double} values (head + tail) for
${\sim}104$~bits of mantissa in 16~bytes.  On GPUs this type is less interesting
than \code{dblflt} because it exercises the slow fp64 units---if fp64
throughput is available, a single \code{double} is usually sufficient for
perturbation, and if more precision is needed, the 4$\times$32 type avoids
fp64 entirely.

\paragraph{4\texorpdfstring{$\times$}{x}32 quad-float (\code{GQF::gqf\_real}).}
Four \code{float} values (${\sim}96$-bit mantissa, ${\sim}29$ decimal digits)
stored as a CUDA \code{float4} vector.  The implementation is adapted from the
QD library by Bailey and Hida \cite{Hida2001QD} (Lawrence Berkeley National Laboratory) and
resides in the \code{GQF} namespace with a \code{make\_qf()} constructor.
Because it uses only fp32 ALUs, it can, in principle, be faster than fp64-based alternatives
on consumer GPUs despite using four components.

\paragraph{4\texorpdfstring{$\times$}{x}64 quad-double (\code{GQD::gqd\_real}).}
Four \code{double} values (${\sim}212$-bit mantissa, ${\sim}64$ decimal
digits) stored as an aligned \code{double4} type.  This type is the highest
software-only precision available before switching to arbitrary-precision
arithmetic (MPIR \cite{MPIR_Library} on CPU or \code{HpSharkFloat} on GPU).  Also adapted from
the QD library.

Unlike \code{dblflt}, the quad types do not have a \code{CudaDblflt}-style
wrapper class; they use direct namespaced helper functions for arithmetic.

\subsection{Precision hierarchy}
\label{subsec:ep-hierarchy}

The numeric roles available in \FractalShark{}, from per-pixel formats through
reference-orbit formats, are:

\begin{center}
\begin{tabularx}{\textwidth}{llrX}
\hline
\textbf{Type} & \textbf{Components} & \textbf{Effective bits} & \textbf{Practical role} \\
\hline
\code{float}                & 1$\times$fp32 & ${\sim}24$  & fastest; sometimes insufficient \\
\code{dblflt}               & 2$\times$fp32 & ${\sim}48$  & practical sweet spot \\
\code{double}               & 1$\times$fp64 & ${\sim}53$  & standard; slow on consumer GPUs \\
\code{dbldbl}               & 2$\times$fp64 & ${\sim}104$ & academic \\
\code{GQF::gqf\_real}       & 4$\times$fp32 & ${\sim}96$  & academic \\
\code{GQD::gqd\_real}       & 4$\times$fp64 & ${\sim}212$ & academic \\
MPIR                         & runtime limbs & runtime-selected & CPU reference recurrence \\
\code{HpSharkFloat}         & fixed 32-bit limbs & fixed by specialization & CUDA reference recurrence \\
\hline
\end{tabularx}
\end{center}

Selected native and expansion mantissa types have \code{HDRFloat}
instantiations when wide dynamic range is required.  In practice, the vast majority of
rendering uses either \code{float}, \code{double}, or
\code{HDRFloat<CudaDblflt<dblflt>>} as the perturbation delta type.

The numeric types introduced here are put to use in the following chapter,
which develops the perturbation framework that makes deep-zoom rendering
practical by reducing per-pixel arithmetic from full precision to lightweight
floating-point deltas.
